\documentclass[aspectratio=169]{beamer}
\usetheme{Madrid}
\usecolortheme{beaver}
\usepackage[utf8]{inputenc}
\usepackage[T1]{fontenc}
\setbeamertemplate{navigation symbols}{}
\setbeamercolor{title}{fg=white,bg=cyan!60!black}
\setbeamercolor{frametitle}{fg=white,bg=cyan!50!black}
\setbeamercolor{block title}{bg=cyan!15,fg=black}
\setbeamercolor{block body}{bg=cyan!5,fg=black}

\title{G1 — Sprite Start}
\subtitle{Coordinates, sequences — draw and move your first sprite}
\author{ByteForge --- Grade 7 --- Game Maker}
\date{}

\begin{document}

\begin{frame}
    \titlepage
\end{frame}

\begin{frame}{Today's Plan}
    \begin{itemize}
        \item First Sprite
        \item Four Corners
        \item Rectangle Room
        \item Three in a Row
        \item[$\bigstar$] BOSS --- Sprite Follows You
    \end{itemize}
\end{frame}

\begin{frame}{Coordinates — every sprite has an (x, y)}
    \Large
    x (how far right), y (how far down) — y grows DOWN, the opposite of a maths graph. Everything you build this year comes back to these two numbers.
\end{frame}


\begin{frame}{Reading coordinates precisely}
    \Large
    "Draw 4 things" and "draw 4 things in specific places" are very different tasks — real challenges almost always care about WHERE, not just how many.
\end{frame}


\begin{frame}{Rectangles need FOUR numbers, not two}
    \Large
    A rectangle needs a position AND a size (width, height) — four numbers instead of a circle's three.
\end{frame}


\begin{frame}{Sequence — same Y, different X}
    \Large
    Three shapes "in a line" means one coordinate stays the SAME while the other changes — a huge amount of real layout is exactly this.
\end{frame}


\begin{frame}{BOSS — coordinates that update every frame}
    \Large
    "On every frame" re-runs its contents dozens of times a second — clear, then redraw, every frame, is the engine behind every animated game.
\end{frame}

\begin{frame}{Real-World Example}
    \Large
    Your phone's touchscreen reporting exactly WHERE you tapped (not just THAT you tapped) — every touch event carries an (x, y), same idea as this chapter's sprite position.
\end{frame}

\begin{frame}{Watch}
    \Large
    Suggested video topic: \emph{Screen coordinates explained — x y position for beginners}

    \vspace{1em}
    \normalsize
    Search this on YouTube together, or swap in a specific video you've
    already vetted --- do not project an unverified link cold in class.
\end{frame}

\begin{frame}{Let Them Discover}
    \Large
    On paper, draw a 480x260 rectangle. Mark where (0,0), (240,130), and (479,259) actually land. Which corner is (0,0) — and why might that surprise someone used to maths graphs?
\end{frame}

\begin{frame}{Your Turn}
    \Large
    Open the student map, find \textbf{G1 — Sprite Start}, and work through it
    hands-on --- everything you just saw, they now do themselves.
\end{frame}

\end{document}
